% Generated draft: requires mathematical and source review.
\documentclass[11pt]{article}
\usepackage[margin=25mm]{geometry}
\usepackage{fontspec}
\setmainfont{Latin Modern Roman}
\usepackage{amsmath,amssymb,mathtools,unicode-math}
\setmathfont{Latin Modern Math}
\usepackage{enumitem,needspace,xurl,hyperref}
\setlength{\parindent}{0pt}
\setlength{\parskip}{5pt}
\setlength{\emergencystretch}{3em}
\newcommand{\N}{\mathbb{N}}
\newcommand{\Z}{\mathbb{Z}}
\newcommand{\Q}{\mathbb{Q}}
\newcommand{\R}{\mathbb{R}}
\newcommand{\C}{\mathbb{C}}
\newcommand{\F}{\mathbb{F}}
\newcommand{\A}{\mathbb{A}}
\newcommand{\PP}{\mathbb{P}}
\newcommand{\E}{\mathbb{E}}
\newcommand{\eps}{\varepsilon}
\newcommand{\dd}{\,\mathrm{d}}
\begin{document}
\section*{UnsolvedMath ID: 1263}
\textbf{Source problem code:} catalogue entry
\subsection{Definitions and mathematical statement}
In the usual nonnegative-integer setting, let $B\subseteq\Z_{\ge0}$ and define the unordered representation function by
\[
 r_B(n)=\bigl|\{(b_1,b_2)\in B^2:b_1\le b_2,\ b_1+b_2=n\}\bigr|,
 \qquad n\in\Z_{\ge0}.
\]
Repetitions $b_1=b_2$ are allowed. The set $B$ is an asymptotic additive basis of order $2$ if
\[
 \exists N_0\in\Z_{\ge0}\ \forall n\ge N_0,\quad r_B(n)\ge1.
\]
The literal question in the supplied record is whether every such $B$ satisfies
\[
 \lim_{n\to\infty}r_B(n)=+\infty,
\]
that is,
\[
 \forall M\in\Z_{\ge1}\ \exists N\in\Z_{\ge0}\
 \forall n\ge N,\quad r_B(n)\ge M.
\]
The domain $B\subseteq\Z_{\ge0}$ is the conventional setting suggested by the named conjecture; the record itself does not explicitly specify $B$'s ambient domain. Allowing arbitrary negative integers is a different problem and is not silently included. The source also leaves ordered versus unordered counting unspecified. If ordered pairs are counted instead, their number $r_B^{\mathrm{ord}}(n)$ obeys
\[
 r_B^{\mathrm{ord}}(n)=2r_B(n)-\delta_B(n),\qquad
 \delta_B(n)=\begin{cases}1&n\text{ even and }n/2\in B,\\0&\text{otherwise},\end{cases}
\]
so the divergence and unboundedness questions are unchanged by those two conventions.

The supplied triage emphasizes that the standard Erd\H{o}s--Tur\'an additive-basis conjecture asks only for unboundedness, equivalently
\[
 \limsup_{n\to\infty}r_B(n)=+\infty,
 \qquad\text{or}\qquad
 \forall M\ge1\ \forall N\ge0\ \exists n\ge N:\ r_B(n)\ge M.
\]
This differs from eventual divergence in the literal stored question; no equivalence is asserted. The background attributes the conjecture to Erd\H{o}s and Tur\'an in 1941 but describes it using the stronger divergence wording. That attribution does not justify replacing either quantifier pattern by the other.

The literature triage checked 17 August 2026 reports that recent literature continues to call the standard unboundedness conjecture unsolved, and that it verified no theorem settling the stronger literal divergence assertion. It therefore asks first for clarification of the target. Its partial-progress classification and the registry's main open label are retained as supplied, rather than treated as a proof of either formulation's status. The background mentions results for special bases without identifying precise results or hypotheses.
\subsection{Short English statement}
Must the number of two-term representations of every sufficiently large integer eventually exceed every fixed bound for any order-two additive basis? The source asks this stronger divergence question, while its named Erd\H{o}s--Tur\'an conjecture traditionally asks only that the counts be unbounded.
\subsection{Sources}
\begin{itemize}
\item Local UnsolvedMath registry, permanent ID $1263$, source code \texttt{catalogue entry}, ``Erd\H{o}s--Tur\'an Conjecture on Additive Bases,'' including the background's 1941 attribution and supplied literature triage checked 17 August 2026.
\item M. B. Nathanson, \emph{Solutions to some problems on unique representation bases}, Journal of Combinatorial Theory, Series A \textbf{221} (2026), 106166, \url{https://doi.org/10.1016/j.jcta.2026.106166}; cited in the record for the standard unbounded-representation-function formulation and its unsolved status.
\end{itemize}
\end{document}
